
\subsubsection{6.1 Interpretation of Results}

The data show that for CodeLlama-7B-Instruct on MBPP runtime errors, minimal feedback (G0, G1) yields higher repair success rates than detailed feedback (G2, G3, G4). This pattern was not predicted by the initial hypotheses.

Several possible explanations exist for this pattern, none of which can be definitively established from the current data:

\begin{enumerate}
\item \textbf{Prompt length}: Detailed feedback increases prompt length, which may affect model attention. G4 prompts can be substantially longer than G0 prompts due to stack trace content.
\end{enumerate}

\begin{enumerate}
\item \textbf{Anchoring effects}: Detailed error information may cause the model to focus on the specific error location even when the root cause lies elsewhere.
\end{enumerate}

\begin{enumerate}
\item \textbf{Model capacity}: At the 7B parameter scale, the model may lack capacity to productively integrate detailed feedback.
\end{enumerate}

These explanations are not mutually exclusive and cannot be distinguished based on the current experiment.

\subsubsection{6.2 Cluster Pattern}

The success rates cluster into two groups separated at the G1-G2 boundary. Within the high-success cluster (G0, G1), the difference is 1.0 percentage point (not significant). Within the low-success cluster (G2, G3, G4), differences range from 1.6 to 5.9 percentage points.

\subsubsection{6.3 Relation to Prior Work}

Self-Debug \citep{chen2023selfdebug} reported 12\% improvement on MBPP using approximately G2-level feedback. The current study's G2 success rate of 18.4\% is in a similar range, though direct comparison is limited by differences in experimental setup. The current study extends prior work by including the G0 condition, which was not systematically evaluated in Self-Debug.

Haque et al. (2025) observed that execution traces provide limited improvement for LLM repair unless prompts are optimized for LLMs. The current results are consistent with this observation.

\subsubsection{6.4 Limitations}

The following limitations constrain the generalizability of these findings:

\textbf{L1: Single Model.} Results are specific to CodeLlama-7B-Instruct. Larger models (13B, 34B, 70B) may show different patterns. No claims are made about optimal granularity for larger models.

\textbf{L2: Single Benchmark.} MBPP consists of relatively simple single-function problems. Results may differ for more complex debugging tasks.

\textbf{L3: Single Prompt Template.} The Self-Debug template was used throughout. Template-granularity interactions are possible but were not tested.

\textbf{L4: Single Repair Attempt.} Only single-turn repair was evaluated. Multi-turn iterative repair may show different patterns.

\textbf{L5: Deterministic Generation.} Temperature 0 was used for reproducibility. Stochastic sampling may yield different results.

\textbf{L6: Runtime Errors Only.} The study focused on runtime errors (60.8\% of failures). Silent logic errors producing wrong outputs may respond differently to granularity variations.

\textbf{L7: Zero Initial Pass Rate.} CodeLlama-7B-Instruct did not successfully solve any MBPP problems in this configuration before repair. This unusual baseline may affect generalizability.

\subsubsection{6.5 Implications}

For CodeLlama-7B-Instruct on MBPP-style problems, minimal error feedback appears more effective than detailed feedback for single-turn repair. Whether this finding transfers to other models, tasks, or configurations requires additional empirical investigation.
